%----------CABECERA----------
\begin{center}
  \textbf{\Huge \scshape Mikel Loinaz Anton}\\[8pt]
  \textbf{Jefe de proyecto/Systems Architect}\\
  \small
  +34 664 502 113
  $|$
  \href{mailto:mikel.loinaz.anton@gmail.com}{\underline{mikel.loinaz.anton@gmail.com}}
  $|$
  Mungia, Bizkaia
\end{center}

%-----------RESUMEN-----------
\section{Resumen Profesional}
\resumeSubHeadingListStart
  \item{\small{
    Ingeniero Aeroespacial con experiencia en gestión de equipos y coordinación de proyectos
    estratégicos en entornos de I+D. Actualmente trabajo como Jefe de Proyecto y Arquitecto de Sistemas. Persona organizada y proactiva con habilidades sociales
    en liderazgo, planificación, gestión del tiempo y resolución de problemas.
  }}
\resumeSubHeadingListEnd

%-----------EXPERIENCIA-----------
\section{Experiencia Laboral}
\resumeSubHeadingListStart

  \resumeSubheading
    {Jefe de Proyecto / Systems Architect}{Mar. 2025 -- Actualidad}
    {GMV A\&D}{Tres Cantos, Madrid}
    \resumeItemListStart
      \resumeItem{Liderazgo técnico de proyectos de I+D: navegación colaborativa para enjambres autónomos heterogéneos, con especial foco en el diseño de comunicaciones robustas entre agentes.}
      \resumeItem{Preparación de ofertas técnicas y coordinación con clientes institucionales (EDA, Comisión Europea), asi como participación en congresos europeos (ASCI, EDA, DESEI+d)}
    \resumeItemListEnd

  \resumeSubheading
    {Ingeniero Aeroespacial}{Mar. 2022 -- Mar. Actualidad}
    {GMV A\&D}{Tres Cantos, Madrid}
    \resumeItemListStart
      \resumeItem{Desarrollo de algoritmos de path planning y task allocation mediante aprendizaje por refuerzo (RL) para navegación autónoma en entornos desconocidos (proyecto AI-GNCAir).}
      \resumeItem{Despliegue de algoritmos de navegación y estimación de estado (Filtro de Kalman Extendido y Error-State Kalman Filter) para sistemas de guiado autónomo.}
      \resumeItem{Generación y etiquetado de datos para entrenamiento de modelos de visión por computador, incluyendo despliegue local de SAM2 y creación de datasets mediante simuladores propios y fuentes externas.}
    \resumeItemListEnd

\resumeSubHeadingListEnd

%-----------FORMACIÓN-----------
\section{Formación}
\resumeSubHeadingListStart
  \resumeSubheading
    {Máster en Sistemas Electrónicos de Información y Comunicación}{Sep. 2026 -- }
    {UNED}{Madrid}
    \resumeSubSubheading{Formación especializada en sistemas de información y telecomunicaciones, electrónica, microelectrónica y comunicaciones inalámbricas, móviles y por satélite}{}
  \resumeSubheading
    {Curso: Análisis de Datos con IA}{Sep. 2024 -- Mar. 2025}
    {Datamecum}{Madrid}
    \resumeSubSubheading{Curso intensivo de 6 meses de IA aplicada al análisis, generación y tratamiento de datos}{}

  \resumeSubheading
    {Grado en Ingeniería Aeroespacial}{Sep. 2017 -- Jun. 2022}
    {Universidad Politécnica de Madrid}{Madrid}
    \resumeSubSubheading{Especialización en Propulsión Aeroespacial}{}
  
\resumeSubHeadingListEnd

%-----------APTITUDES E IDIOMAS-----------
\section{Aptitudes Técnicas e Idiomas}
\begin{itemize}[leftmargin=0.15in, label={}]
  \small{\item{
    \textbf{Herramientas}{: Enterprise Architect, Jira, Office} \\
    \textbf{Programación}{: C++, Python, Docker, PyTorch} \\
    \textbf{Idiomas}{: Español (nativo), Euskera (nativo, C2), Inglés (C1)}
  }}
\end{itemize}
